// While Loop — Newton's method for square root
//
// Educational example demonstrating while loops with a per-pixel stop.
// Uses Newton-Raphson iteration to compute the square root of the
// pixel luminance: guess = (guess + x / guess) / 2
// Iterates until the error is below a tolerance, then compares the
// result with the built-in sqrt() function.
//
// The error differs from pixel to pixel, so a `break` on it would end the
// loop for EVERY pixel on the first pass (LANGUAGE.md §7.1). Each pixel keeps
// a `done` flag instead and stops updating its guess once it has converged,
// while the loop itself runs $max_iter passes, the same for every pixel.
//
// Output: left half shows Newton's result, right half shows the
// difference (error) magnified 100x. Should be near-black.

f$tolerance = 0.0001;  // convergence threshold
i$max_iter = 20;        // maximum iterations to prevent infinite loops

vec3 src = @image;
float x = luma(src);

// Newton's method for sqrt(x)
float guess = max(x, 0.001);   // initial guess (avoid zero)
float done = 0.0;              // 1.0 once this pixel has converged
int iter = 0;

while (iter < $max_iter) {
    float next = (guess + x / max(guess, 0.0001)) * 0.5;
    float err = abs(next - guess);
    if (done < 0.5) {
        guess = next;          // a converged pixel keeps its guess
    }
    if (err < $tolerance) {
        done = 1.0;
    }
    iter = iter + 1;
}

float newton_sqrt = guess;
float builtin_sqrt = sqrt(x);

// Visualize: left = Newton result, right = error * 100
if (u < 0.5) {
    @OUT = vec3(newton_sqrt);
} else {
    float diff = abs(newton_sqrt - builtin_sqrt) * 100.0;
    @OUT = vec3(diff, 0.0, 0.0);
}
